% ---------------------------------------------------------------------------
% Packages
% ---------------------------------------------------------------------------
% This template targets pdfLaTeX because it is widely available on Overleaf,
% TeX Live, and MiKTeX. If you switch to LuaLaTeX or XeLaTeX, remove inputenc.

\usepackage[utf8]{inputenc}
\usepackage[T1]{fontenc}
\usepackage[ngerman]{babel}
\usepackage{lmodern}
\usepackage{csquotes}
\usepackage{etoolbox}

\usepackage[a4paper,left=3cm,right=2.5cm,top=2.5cm,bottom=2.5cm,headheight=14pt]{geometry}
\usepackage{setspace}
\usepackage{microtype}

\usepackage{xcolor}
\usepackage{graphicx}
\usepackage{booktabs}
\usepackage{tabularx}
\usepackage{enumitem}
\usepackage{amsmath,amssymb,amsthm}
\usepackage{siunitx}

\usepackage{upquote}
\usepackage{listings}
\usepackage[font=small,labelfont=bf]{caption}
\usepackage{subcaption}
\usepackage[nohyperlinks]{acronym}
\usepackage{fancyhdr}
\usepackage{titlesec}
\usepackage[nottoc]{tocbibind}
\usepackage[
  colorlinks=true,
  linkcolor=black,
  citecolor=black,
  urlcolor=blue,
  pdfusetitle
]{hyperref}
\usepackage{bookmark}

% ---------------------------------------------------------------------------
% PDF metadata
% ---------------------------------------------------------------------------
\hypersetup{
  pdftitle={\thesistitle},
  pdfauthor={\thesisauthor},
  pdfsubject={\thesistype},
  pdfkeywords={\pdfkeywords}
}

% ---------------------------------------------------------------------------
% Layout
% ---------------------------------------------------------------------------
\onehalfspacing
\setlength{\parindent}{0pt}
\setlength{\parskip}{1.0ex}

% Keep the table of contents readable: chapter, section, and subsection are
% usually enough for a thesis outline.
\setcounter{tocdepth}{2}
\setcounter{secnumdepth}{2}

\sisetup{
  locale=DE,
  per-mode=symbol
}

\pagestyle{fancy}
\fancyhf{}
\fancyhead[L]{\small\nouppercase{\leftmark}}
\fancyhead[R]{\small\thepage}
\renewcommand{\headrulewidth}{0.4pt}

\titleformat{\chapter}[hang]
  {\normalfont\huge\bfseries}
  {\thechapter}
  {1em}
  {}
\titleformat{name=\chapter,numberless}[hang]
  {\normalfont\huge\bfseries}
  {}
  {0pt}
  {}
\titlespacing*{\chapter}{0pt}{0pt}{2.5em}

% ---------------------------------------------------------------------------
% Source-code listings
% ---------------------------------------------------------------------------
% Use language=Python, Java, C++, SQL, HTML, XML, etc. as needed.
% Long listings are usually better placed in the appendix.
\definecolor{listingbackground}{HTML}{F8F6FC}
\definecolor{listingrule}{HTML}{5940D9}
\definecolor{listingnumber}{HTML}{8A8494}
\definecolor{listingkeyword}{HTML}{5940D9}
\definecolor{listingcomment}{HTML}{6F6A75}
\definecolor{listingstring}{HTML}{982D5D}

\lstdefinestyle{thesislisting}{
  basicstyle=\ttfamily\small,
  backgroundcolor=\color{listingbackground},
  keywordstyle=\bfseries\color{listingkeyword},
  commentstyle=\itshape\color{listingcomment},
  stringstyle=\color{listingstring},
  numbers=left,
  numberstyle=\scriptsize\color{listingnumber},
  stepnumber=1,
  numbersep=14pt,
  frame=leftline,
  rulecolor=\color{listingrule},
  framesep=8pt,
  framerule=2pt,
  xleftmargin=1.6em,
  breaklines=true,
  breakatwhitespace=true,
  columns=fullflexible,
  keepspaces=true,
  tabsize=2,
  captionpos=b,
  showstringspaces=false,
  upquote=true
}
\lstset{style=thesislisting}
\renewcommand{\lstlistlistingname}{Quellcodeverzeichnis}

% ---------------------------------------------------------------------------
% Helper commands
% ---------------------------------------------------------------------------
\newcommand{\code}[1]{\texttt{#1}}
\newcommand{\placeholder}[1]{\textcolor{red!70!black}{#1}}
